% commentaire.tex — f169r, vv. 551–589 (conversion initiale des travaux antérieurs)
\comm{556–7}{ὁρᾷς ὡς τλεγγίδας λαβοῦσαι / ἔπειτα σιφωνίζομεν τὸν σῖτον.}{Pour le texte (τὸν κύσθον pour τὸν σῖτον de R) et le sens de σιφωνίζω, voir le commentaire critique. Profitons-en pour remarquer à quel point les griefs des femmes envers Euripide ne sont pas sérieux dans la bouche d'Aristophane (si Euripide est misogyne, c'est bien là le moindre de ses défauts) ; car si les maris athéniens, s'en revenant d'une pièce d'Euripide, regardaient d'un œil soupçonneux leur femme tenant une marmite, avec quels yeux vont-ils dorénavant observer leur femme se saisir d'un strigile ou de tout instrument de forme analogue ?! Quoi qu'il en soit, comme pour l'histoire de l'amant qui s'échappe (v. 501), Mnésiloque n'invente rien mais reprend des thèmes certainement populaires du Landerneau athénien : le détournement à fin lubrique du strigile se retrouve illustré, ce me semble, sur une coupe d'Onésimos, antérieure à la pièce d'Aristophane. Pour finir – était-il besoin de préciser ?.. – précisons tout-de-même : Mnésiloque mime en lançant son « ὁρᾷς » à Mika : « Ainsi font, font, font, les petites midinettes. Un siphon, fond, fond, on se siphonne toutes à fond.» \nc{Bâle, Antikenmuseum BS 440. ARV 326, 86 \italique{bis}. \italique{Para} 359. Elle est reproduite par John Boardman, fig. 233 in \italique{Les vases athéniens à figures rouges}, \italique{La période archaïque} ; première édition sous le titre \italique{Athenian Red Figure Vases}, \italique{The Archaic Period}, 1975. (Le lecteur pourra la retrouver aux annexes de ce livre.) En parcourant le livre de Boardman, cette coupe a retenu mon attention : elle me paraît parfaitement illustrer l'anecdote de Mnésiloque. (Reconnaissons tout d'abord que l'interprétation est sujette à caution : j'ai montré cette coupe à différents amis et ai eu droit à autant d'interprétations qu'il y avait d'amis... interprétations qui en disaient d'ailleurs plus long sur leur propre psychologie que sur les mœurs des Athéniens ; profitons-en donc pour suggérer à l'éventuel psychanalyste qui se serait égaré à lire ces lignes d'abandonner fissa ses infantiles tests de Rorschach et de passer aux coupes athéniennes du \nieme{VI} siècle : dix vases athéniens d'Onésime seraient aussi révélateurs des obsessions ou perversions des supposés psychotiques que les taches d'encre du regretté Hermann, arraché trop tôt à l'affection des siens... Sans parler du gain en valeur artistique : incommensurable ! Refermons cette parenthèse et retournons à notre coupe :) qu'y ai-je donc vu ? À droite, un satyre (sa calvitie naissante est caractéristique des représentations de satyres de cette époque ; c'est entre autre sur ce détail de la laideur physique de Socrate, visible sur ses bustes, qu'Alcibiade rapproche selon moi Socrate du satyre Silène), un satyre donc, en train de se strigiler le fondement, pointe (du doigt) le Mont de Vénus d'une femme nue, lui désignant par là l'usage impropre du strigile auquel il lui arrive de céder. La clef qui me manque pour être tout-à-fait sûr de la bonne compréhension de cette scène est l'identification de l'objet dans la main de la femme : je crois y reconnaître un strigile enroulé dans du tissu. (plus recommandé pour atteindre le plaisir que la lame de bronze brute.) D'aucun y a vu le pagne du jeune homme à gauche. J'interprète cette scène comme un résumé des vices féminins de l'imaginaire grec : la femme est incapable de se contenir, elle cède à tous les plaisirs : le plaisir de la vigne (la femme est adossée à une outre) ; le plaisir de la chair (l'adultère est suggéré avec l'amant à gauche qui se rhabille) et -- toujours insatiable ! -- le plaisir des dieux. (nue, la voilà habillée pour l'hiver !..)
Une autre hypothèse d'un mien ami quant à cette scène vaut d'être rapportée : il l'avait interprétée quelque peu différemment : tout en restant dans le domaine d'Onan, il proposait de garder τὸν σῖτον et de comprendre : \italique{le pain} [d'argile] servant à récupérer les peaux mortes \& autres hormones raclées avec l'huile à l'aide du strigile (les στλεγγίσματα et autre γλοιός), raclures qui étaient ensuite parfois monnayées (notamment lorsqu'elles provenaient de virils athlètes)... et pourquoi pas par quelque vicelarde pour s'en foutre plein les zones érogènes ? Son explication avait l'avantage indéniable de conserver la leçon du manuscrit... Mais je crains que son imagination ne se soit quelque peu emballée : je n'ai pas trouvé d'exemple d'acception du mot σῖτος pour désigner ce «~pain d'argile~». Et si Pline nous confirme un usage médical des στλεγγίσματα, l'usage détourné, qui lui a été suggéré par certains papiers de « recherche », est un pur fantasme des modernes scholiastes.}}
\comm{558–9}{ὥς τ᾽ αὐτὰ κρέ᾽ ἐξ Ἀπατουρίων ταῖς μαστροποῖς διδοῦσαι / ἔπειτα τὴν γαλῆν φαμεν.}{Si Mnésiloque avait commencé un discours tout en nuance, il a bien vite perdu de sa prudence et se met à surpasser le maître Euripide dans la diffamation du genre féminin. Le ton monte et Mnésiloque en est maintenant à dénoncer les pratiques les plus honteuses des Athéniennes les plus dissolues.\\
Nous ne nous attarderons pas sur ce que pouvaient bien procurer les procureuses en échange de ces détournements de viandes\nc{Contentons-nous de renvoyer à la magistrale note d'Alphonse Willems sur la fête des Apaturies, t. II, pp. 118-120 ; le lecteur y découvrira l'importance démocratique de la κρεανομία, l'institution de la distribution de viande. Sur les débordements que de telles orgies pouvaient entraîner, l'on pourra – un exemple parmi d'autres – essayer de comprendre le jeu de mots d'Aristophane sur l'ἀνάρρυσις (au vers 890 de la \italique{Paix}), le second jour des Apaturies.}.
Plus intéressante est l'excuse alléguée pour expliquer cette soudaine disparition carnée. La \italique{team 1\expo{er}} degré des commentateurs d'Aristophane n'a pas manqué de faire remarquer que la belette était un animal chapardeur\nc{Voir par exemple les vers 363-364 des \italique{Guêpes}. Lucien, \italique{Piscator}, 34 : ἁρπακτικώτεροι τῶν γαλῶν.}. Certes ! mais ce pauvre animal a surtout \italique{bon dos} : on l'accuse bien facilement ; ainsi aux vers 1150-1152 de la \italique{Paix} :\citComm{ἦν δὲ καὶ πυός τις ἔνδον καὶ λαγῷα τέτταρα,\\
εἴ τι μὴ 'ξήνεγκεν αὐτῶν ἡ γαλῆ τῆς ἑσπέρας·\\
ἐψόφει γοῦν ἔνδον οὐκ οἶδ᾽ ἅττα κἀκυδοιδόπα\\
Il y avait aussi au logis de l'amouille et quatre morceaux de lièvre ; 
à moins que la belette n'en ait volé hier soir ; 
car il se passait là-dedans je ne sais quel bruit et quel remuement.}L'allusion au remue-ménage n'est pas limpide mais il est clair que le sens ne peut être pris littéralement : qui, entendant une belette le dévaliser, ne s'armerait aussitôt d'un bâton pour châtier la voleuse ? (Noter d'ailleurs ici aussi la présence de l'article : ἡ γαλῆ.) 
Aristophane joue certainement sur un détournement d'une locution proverbiale : la belette attique étant notre \italique{bouc-émissaire} à qui l'on fait \italique{porter le chapeau}. 
Conformité du grec et du français : la belette joue le même rôle que notre chat, animal qui n'a pas encore été introduit dans l'Athènes du \nieme{V} siècle : il chasse les rongeurs. Or, que lit-on chez La Fontaine ?\citComm{Le sommeil quitta son logis, /
Il eut pour hôte les soucis, /
Les soupçons, les alarmes vaines. / 
Tout le jour il avait l'œil au guet ; \textbf{et la nuit, /
Si quelque chat faisait du bruit, /
Le chat prenait l'argent.}\nc{Il s'agit – vous l'aurez bien sûr reconnue ; sinon, cessez immédiatement votre lecture et ne revenez que lorsque vous la connaîtrez par cœur ! – de la seconde fable du huitième livre : \italique{Le Savetier et le Financier}.}}Ce n'est pas moi qui m'abuse : remplacez l'argent par la nourriture et les deux passages disent exactement la même chose.
Voici ce que l'on trouve dans le \italique{Dictionnaire de la langue verte} d'Alfred Delvau :\citComm{« C'est le chat ! » : Expression de l'argot du peuple, qui souligne ironiquement un doute, une dénégation. Ainsi, quelqu'un disant : Ce n'est pas moi qui ai fait cela. – Non ! c'est le chat, lui répondra-t-on.}Un autre passage de la \italique{Paix} montre encore mieux ce sens imagé : aux vers 793-795 :\citComm{καὶ γὰρ ἔφασχ᾽ ὁ πατὴρ ὃ παρ᾽ ἐλπίδας /
εἶχε τὸ δρᾶμα γαλῆν τῆς /
ἑσπέρας ἀπάγξαι.\\
Aussi bien leur père prétend-il qu'un drame qu'il pensait tenir contre toute attente fut étranglé le soir par une belette.}Juste avant ce passage, Aristophane faisait un honnête résumé des bienfaits qu'il a apportés à la Comédie, il en vient ensuite à étriller les nullités de ses poétereaux de contemporains. Nous en sommes à Carcinos \& Fils qui font des pieds et des mains pour pouvoir danser – c'est-à-dire obtenir la subvention d'un chœur.
Mais une fois le chœur accordé – à rebours de ce qui était attendu ! (παρ᾽ ἐλπίδας) [du bon goût ?] – Carcinos-père fait volte-face et prétend que son drame « fut étranglé le soir par une belette » (c'est encore le soir que la belette commet son forfait). 
Le scholiaste nous est ici précieux puisqu'il précise que Carcinos voulait représenter une pièce intitulée les \italique{Souris} (Μῦες\nc{Probablement le drame satyrique. Μῦς étant un nom propre, comme chez Hérodote (VIII, 133), dans la \italique{Souda} (ο 680) ou chez Athénée (XI, 782b), les Μῦες désignant les enfants de ce Μῦς. Je m'inscris en faux contre l'interprétation de Paul Mazon (\italique{La} Paix\italique{, texte grec publié avec une introduction, des notes critiques et explicatives} (1904), v. 793) reprise par Taillardat (§792, p.463 n.5) qui pensent que Μῦες était le surnom de ces débauchés de Carcinites. L'on s'expliquerait difficilement la confusion du scholiaste (présent sur tous les bons manuscrits et repris par Suidas κ 396). Objection plus forte : cela n'expliquerait pas le texte de la \italique{Paix} : la belette (γαλῆ) est accusé d'avoir étranglé nuitamment non les choreutes mais la pièce elle-même (τὸ δρᾶμα).}) ; encore que la chronique du temps ne nous ait pas conservé l'affaire, il paraît aisé de reconstituer peu ou prou ce qui s'est passé. Les poètes devaient présenter à l'archonte éponyme leur projet. 
Pour les Dionysies, l'archonte sélectionnait trois tétralogies qui recevaient chacune le soutien financier d'un riche Athénien : le chorège. 
Ce n'est qu'alors que le poète se mettait à écrire. Et c'est à cette étape que les forces auront probablement manqué à notre Carcinos… ce dernier croyant s'en tirer avec un trait d'esprit : s'il ne peut pas présenter ses \italique{Souris}, c'est que la belette les a étranglées. 
Le mot ne peut être \italique{bon} que s'il existait effectivement une expression pour faire endosser la responsabilité d'un méfait à cette belette proverbiale.}
\comm{568}{λαβὲ θοἰμάτιον Φιλίστη.}{Mettre bas les manteaux était de mise avant d'en venir aux mains. Ainsi au vers 408 des \italique{Guêpes} ; autre exemple : le fragment d'Hipponax\nc{Fr. 121 de l'édition Degani : \italique{Hipponactis testimonia et fragmenta}, Teubner, 1991 ; fr. 120 chez Olivier Masson, voir aussi son introduction p. 13.} contre Boupalos\nc{Il est question de ce Boupalos au vers 361 de \italique{Lysistrata} -- le fragment sus-cité d'Hipponax nous a d'ailleurs été transmis par le scholiaste à ce passage ; je reprends la note de Willems de ce passage, p. 351 du second tome.} : Boupalos, sculpteur de Chios, avait fait d'Hipponax un portrait où il exagérait encore la laideur du modèle. Pour se venger, le poète écrivit contre lui des vers mordants et pleins de menaces. Suidas (β 452) nous a conservé celui-ci, qui paraît avoir passé en proverbe : \citComm{λάβετε μευ ταἰμάτια [θοιμάτιον codd.], κόψω Βουπάλου τον οφθαλμόν.\\
\francais{Prenez mon manteau, que je fasse sauter un œil à Boupalos.}}}
\comm{569}{πρόσθες μόνον.}{Willems substitue πρόσθιγε (du verbe προσθιγγάνω, \italique{toucher}) à πρόσθες qu'il soupçonne de ne pas être le mot propre. Il s'en explique succinctement dans la \nieme{16} note reproduite à la fin. D'aucuns ont essayé de justifier la leçon du texte en avançant une tournure elliptique : πρόσθες [τὴν χεῖρα] (expression qui se trouve effectivement au vers 1669 des \italique{Phéniciennes} d'Euripide\nc{L'on cite d'autres exemples équivalents : le vers 933 des \italique{Nuées}: κλαύσει, τὴν χεῖρ᾽ ἢν ἐπιβάλλῃς ; le fragment 309(277) (éd. K.-A. = \italique{Zonarae lexikon} p.1168) : τὴν χεῖρα μὴ ᾿πίβαλλε, μὴ κλάων καθῇ ; \italique{Lysistrata}, vers 435-436 \& 439-440 : τὴν χεῖρά μοι ... προσοίσει ... κλαύσεται \& εἴ ... μόνον τὴν χεῖρ᾽ ἐπιβαλεῖς, ἐπιχεσεῖ πατούμενος.}) mais tous les exemples proposés comportent χεῖρα, qui semble indispensable au sens. Cobet avait donc proposé de lire πρόσιθι μόνον, \italique{approche seulement} ; ce que l'on dirait effectivement dans nos langues modernes mais le grec use d'une image différente : il dit \italique{jette(-moi) seulement ta main} (cf. les exemples déjà cités) ou bien \italique{touche(-moi) seulement} (cf. les exemples proposés par Willems auxquels on peut ajouter \italique{Lysistrata} v. 365, qui est l'exacte correspondance de son émendation : ἅψαι μόνον).}
\comm{570}{τὸν σησαμοῦνθ᾽ ὃν κατέφαγες, τοῦτον χεσεῖν ποιήσω.}{«~Le gâteau de sésame que tu as engoulé, je vais te le faire chier.~» Mnésiloque aurait pu ajouter : que tu viens d'engouler \italique{en cachette} ; car ce \grec{σησαμοῦς}\nc{Suidas (π 3192) nous apprend que c'est une sorte de gâteau fait à partir de miel bouilli et de sésame : Πυραμοῦς: εἶδος πλακοῦντος ἐκ μέλιτος ἑφθοῦ καὶ πυροῦ πεφρυγμένου, ὡς σησαμοῦς ὁ διὰ σησαμῆς. ταῦτα δὲ ἐτίθεσαν ἆθλα τοῖς διαγρυπνηταῖς ἐν τοῖς συμποσίοις, καὶ ὁ διαγρυπνήσας μέχρι τῆς ἔω ταῦτα ἐλάμβανεν ἔπαθλον. ἡμέτερος ὁ πυραμοῦς.} a été avalé en plein milieu des Thesmophories, le jour du \italique{jeûne}, \grec{νηστεία}.
Athénée (XIV, 647a) nous rapporte que les Syracusaines préparaient des gâteaux semblables de miel et de sésame, les \italique{mylles} pour leurs Thesmophories~:\citComm{ΜΥΛΛΟΙ. Ἡρακλείδης ὁ Συρακόσιος ἐν τῷ περὶ Θεσμῶν ἐν Συρακούσαις φησὶ τοῖς Παντελείοις τῶν Θεσμοφορίων ἐκ σησάμου καὶ μέλιτος κατασκευάζεσθαι ἐφήβαια γυναικεῖα, ἃ καλεῖσθαι κατὰ πᾶσαν Σικελίαν μυλλοὺς καὶ περιφέρεσθαι ταῖς θεαῖς.\\
\francais{Mylloi. Héraclide de Syracuse rapporte ce qui suit, dans son \italique{Traité des Rites religieux} : « Le jour des grandes Thesmophories, on fait à Syracuse avec de la farine, du miel et du sésame, la figure de la partie sexuelle d'une femme, que l'on porte en pompe, pour l'offrir aux déesses. On l'appelle \italique{mylle} par toute la Sicile. »}}C'est de peur que Mica risque de relâcher son gâteau. Chez Aristophane la peur se manifeste presque toujours physiquement par des troubles du bas-ventre\nc{On en ramasse à la pelle de ces exemples de déjections consécutives à une frayeur ; Taillardat (§309) en cite bon nombre : \italique{Assemblée des femmes}, v. 1060 ; \italique{Acharniens}, v. 350 ; et d'autres, mais la liste n'est toujours pas exhaustive : on pourrait encore la compléter avec les \italique{Guêpes}, v. 627 ; \italique{Oiseaux}, v. 66 ; etc. etc. Signalons la même promesse de relâchement des intestins après correction dans \italique{Lysistrata}, v. 440 : \grec{ἐπιχεσεῖ πατούμενος}.}. On dira en français \italique{avoir les pétoches}. \latin{Pulchre, bene, recte.} Mais l'essentiel nous échappe encore : Mnésiloque ne promet pas à Mica un petit coup en cas d'agression de sa part, mais une riposte non graduée : une pluie de horions à lui assainir le côlon ! Car il faut comprendre : «~tu vas en chier [de peur] \italique{malgré} le gâteau de sésame par toi engueulé~».
Ce genre de gâteaux avait la vertu d'être astringent\nc{Hypothèse, certes corroborée par les textes d'Aristophane que nous allons voir mais point testée sur mon propre appareil digestif. Est-ce digne d'un disciple de Willems, ce grand amateur de \italique{Realia} ?}. C'est ce qu'il ressort d'un passage des \italique{Nuées} (vers 506-508) : Strepsiade hésite à rentrer dans l'antre de Socrate, qu'il compare aimablement à celui de Trophonius : \citComm{\ \ \ \ \ \ \ \ \ \ \ \ \ \ \ \ \ \ \ \ \ \ \ \ \ \ \ \ \ \ εἰς τὼ χεῖρέ νυν\\
δός μοι μελιτοῦτταν πρότερον, ὡς δέδοικ᾽ ἐγὼ\\
εἴσω καταβαίνων ὥσπερ εἰς Τροφωνίου.\\
Mets-moi d'abord en main un gâteau au miel ; car j'ai peur de descendre, comme si je pénétrais dans l'antre de Trophonios.}Il faut bien-sûr prononcer d'un trait `εἰς Τροφωνίου' pour faire entendre στρόφος, la \italique{colique}\nc{En passant, Euripide use du même jeu de mots scatologique : au vers 300 de son \italique{Ion} : σὺν ἀνδρί $[$Créuse dit à Ion qu'elle est venue avec son mari, mais... ] σηκοὺς δ᾽ ἐκστρέφει Τροφωνίου... que le mari, Xuthus, était d'ailleurs véritablement allé consulter. Le tour était proverbial et Érasme l'a recensé (\italique{Adages} n° 677)}.\\
C'est aussi ce qui explique le bon mot de table qu'Athénée (I, 4d) rapporte :\citComm{ἀμίδα δότω τις ἢ πλακοῦντά τις δότω\\
Qu'on m'apporte un pot de chambre ou bien un gâteau !}Au contraire de tous ces gâteaux de miel, la συρμαία (sorte de rave noire que prisaient les Égyptiens pour leur régime ; voir le commentaire \italique{infra} aux vers 855-857) possède des vertus purgatives. 
Cette opposition est vraisemblablement la raison de la confrontation dans la pièce perdue des \italique{Danaïdes}\nc{Pollux (VII 199 $=$ fr. 276(265) Kassel-Austin). Cf. aussi le passage -- corrompu -- d'Athénée (XIV, 645e) où il est question de l'ἐγκριδοπώλης dans les \italique{Danaïdes}.} entre deux vendeurs : celui de syrméa et celui d'ἐγκρίδες\nc{Athénée (XIV, 645e) : ΕΓΚΡΙΔΕΣ. πεμμάτιον ἑψόμενον ἐν ἐλαίῳ καὶ μετὰ τοῦτο μελιτούμενον.} :\citComm{ἐν δ᾿ Ἀριστοφάνης Δαναΐσι συρμαιοπῶλαι καὶ ἐγκριδοπῶλαι.}}
\comm{575–6}{}{Sur la conservation de ces deux vers que d'aucuns auraient voulu biffer, voir la note de Willems n° 17.}
\comm{578}{ὀλίγῳ τι πρότερον κατ᾽ ἀγορὰν λαλούμενον.}{Clisthène rapporte les cancans de l'Agora, qui font grand cas de ce beau-père d'Euripide introduit chez les femmes pou plaider sa cause. Ce n'est certainement pas Euripide qui est à l'origine de la fuite et ce pour des raisons évidentes : \ \ \ 1° Si Aristophane, l'avait imaginé ainsi, il nous l'aurait fait comprendre directement ou par quelque allusion \ \ \ 2° Euripide n'a aucun intérêt à voir son complice être démasqué, qui ne pourra donc plus plaider sa cause \ \ \ 3° Euripide a beaucoup de défauts, mais pas celui d'être ingrat au point d'envoyer à une mort certaine son dévoué parent. \ \ \ Savoir quel domestique d'Agathon, quel voisin, ayant aperçu Mnésiloque attifé en femme, sera allé colporter la nouvelle sur l'Agora... voilà bien des détails sans importance sur lesquels Aristophane n'a pas daigné s'arrêter ; nous aurions fait de même si les professeurs parmi les plus admis ne s'étaient lancés dans ce genre d'extravagantes et vaines conjectures.}
\comm{579–87}{ἥκω φράσων...}{Passées les retrouvailles, Clisthène prend un ton paratragique -- l'heure est grave ! Noter à ce propos les redoublements\nc{Redoublements propres au style eschyléen et raillés par l'Euripide des \italique{Grenouilles} (vv. 1153 sqq.).} : φράσων/ἀγγελῶν, σκοπῆτε/τηρῆτε, δεινὸν/μέγα, βουλεύοισθε/μέλλοιτε.}
\comm{582–3}{τί δ᾽ ἔστιν ὦ παῖ; παῖδα γάρ σ᾽ εἰκὸς καλεῖν, /	ἕως ἂν οὕτως τὰς γνάθους ψιλὰς ἔχῃς.}{Pique de la coryphée analogue à celle du chœur des \italique{Guêpes}, aux vers 1297-1298 : \citComm{Χορός – τί δ᾽ ἔστιν ὦ παῖ; παῖδα γάρ, κἂν ᾖ γέρων,\\
καλεῖν δίκαιον ὅστις ἂν πληγὰς λάβῃ.}}
\comm{585}{ἀναπέμψαι.}{Le suffixe `ἀνα' s'explique par le fait que le Thesmophorion se trouvait probablement sur la colline de la Pnyx. (Cf. le vers 658.)}
